\subsection{代数的导子}

设 A 为分次 K-代数；对每个齐次元素 \(a \in \mathbf{A}\) ，用 \(\operatorname{ad}_{\varepsilon}(a)\) （若不致混淆，简记为 \(\operatorname{ad}(a)\) ）表示 A 到 A 的 K-线性映射

\[
x \mapsto [ a, x ] _ {\varepsilon}
\]

（第 4 号，公式 (10)），它是次数为 deg a 的分次映射。

命题 8. 设 A 为分次 K-代数。

\begin{enumerate}
\def\labelenumi{(\roman{enumi})}
\tightlist
\item
  对每个 \(\varepsilon\) -导子 \(d: \mathbf{A} \to \mathbf{A}\) 以及每个齐次元素 \(a\) （ \(\mathbf{A}\) 中的），
\end{enumerate}

\[
[ d, \mathrm{ad} _ {\varepsilon} (a) ] _ {\varepsilon} = \mathrm{ad} _ {\varepsilon} (d a).\tag{24}
\]

\begin{enumerate}
\def\labelenumi{(\roman{enumi})}
\setcounter{enumi}{1}
\item
  若代数 \(\mathbf{A}\) 是结合的，则 \(\operatorname{ad}_{\varepsilon}(a)\) 是一个 \(\varepsilon\) -导子（ \(\mathbf{A}\) 的，次数为 \(\deg(a)\) ）。
\item
  假设 \(d\) 的次数为 \(\delta\) ，设 \(\alpha = \deg a\) ，并记 \(f = [d, \operatorname{ad}_{\varepsilon}(a)]_{\varepsilon}\) 。对每个齐次元素 \(x \in A\) （次数为 \(\xi\) ），由 (1)（第 1 号），我们有
\end{enumerate}

\[
\begin{array}{r l} f (x) & = d (a x - \varepsilon (\alpha , \xi) x a) - \varepsilon_ {\delta , \alpha} (a (d x) - \varepsilon_ {\alpha , \delta + \xi} (d x) a) \\ & = (d a) x + \varepsilon_ {\delta , \alpha} a (d x) - \varepsilon_ {\alpha , \xi} (d x) a - \varepsilon_ {\delta + \alpha , \xi} x (d a) \\ & \quad - \varepsilon_ {\delta , \alpha} a (d x) + \varepsilon_ {\alpha , \xi} (d x) a \\ & = (d a) x - \varepsilon_ {\delta + \alpha , \xi} x (d a) = [ d a, x ] _ {\varepsilon}. \end{array}
\]

\begin{enumerate}
\def\labelenumi{(\roman{enumi})}
\setcounter{enumi}{1}
\tightlist
\item
  对所有齐次的 \(x\) （次数为 \(\xi\) ）以及所有齐次的 \(y\) （次数为 \(\eta\) ，在 \(\mathbf{A}\) 中），
\end{enumerate}

\[
\begin{array}{r l} \mathrm{ad} _ {\varepsilon} (a) (x y) & = a (x y) - \varepsilon_ {\alpha , \xi + \eta} (x y) a \\ & = (a x - \varepsilon_ {\alpha , \xi} x a) y + \varepsilon_ {\alpha , \xi} x (a y - \varepsilon_ {\alpha , \eta} y a) \\ & = \mathrm{ad} _ {\varepsilon} (a) (x). y + \varepsilon_ {\alpha , \xi} x. \mathrm{ad} _ {\varepsilon} (a) (y) \end{array}
\]

这里用到了 (1) 以及 A 的结合性。

当 A 是结合的时， \(\operatorname{ad}_{\varepsilon}(a)\) 被称为由 a 定义的 A 的内 \(\varepsilon\) -导子。

推论。设 A 是一个结合的分次代数。对于 A 的两个齐次元素 a, b，

\[
[ \mathrm{ad} _ {\varepsilon} (a), \mathrm{ad} _ {\varepsilon} (b) ] _ {\varepsilon} = \mathrm{ad} _ {\varepsilon} ([ a, b ] _ {\varepsilon}).\tag{25}
\]

只需在 (24) 中把 \(d\) 换成 \(\operatorname{ad}_{\varepsilon}(a)\) ，并把 \(\operatorname{ad}_{\varepsilon}(a)\) 换成 \(\operatorname{ad}_{\varepsilon}(b)\) 。

若 \(\deg a = \alpha\) , \(\deg b = \beta\) ，则公式 (25) 等价于下述关系：对每个齐次元素 \(c \in \mathbf{A}\) （次数为 \(\gamma\) ），

\[
\varepsilon_ {\alpha , \gamma} [ a, [ b, c ] _ {\varepsilon} ] _ {\varepsilon} + \varepsilon_ {\beta , \alpha} [ b, [ c, a ] _ {\varepsilon} ] _ {\varepsilon} + \varepsilon_ {\gamma , \beta} [ c, [ a, b ] _ {\varepsilon} ] _ {\varepsilon} = 0\tag{26}
\]

称为雅可比恒等式。

